\chapter{Hexagonal Architecture}
\label{ch:hexagonal}

Hexagonal Architecture, also known as \textit{Ports and Adapters}, was introduced by Alistair Cockburn to protect application behavior from infrastructure concerns.\autocite{cockburn2005hexagonal} The key idea is simple: we keep application logic in the center, and let external technologies connect through explicit boundaries. These boundaries are called \textit{ports}, and are generally just interfaces in your code. A modern introduction with practical examples is provided by Huttunen.\autocite{huttunen2026hexagonal}

\section{Why We Use Hexagonal Architecture}
\label{sec:hexagonal_why}

Many systems start in a layered style, but business rules often leak into controllers, ORMs, and SDK-specific code. Basically just places where it does not belong. This coupling makes change expensive: replacing persistence technology, adding a new interface, or testing behavior can require touching several technical layers.

Similarly to a 3-layered architecture, we want to separate the concerns of the application into different parts.

Hexagonal architecture addresses this by separating concerns into an \textit{inside} and an \textit{outside}. The inside contains the behavior we care about, while the outside contains delivery and infrastructure details. Comparing to the 3-layered architecture, the presentation and persistence will be considered the "outside" of the application.

\begin{figure}[htbp]
    \centering
    \begin{tikzpicture}[
    core/.style={regular polygon, regular polygon sides=6, shape border rotate=30, draw=blue!70, fill=blue!8, very thick, minimum size=5.4cm, align=center},
    note/.style={align=center, font=\small}
]
    \node[core] (core) {Application\\Core};
    \node[note] at (0,3.4) {Business behavior is protected\\from external technologies};
\end{tikzpicture}

    \caption{The first step is to define a protected application core.}
    \label{fig:hexagonal_core_only}
\end{figure}

\subsection{Purpose of the Inner Hexagon}
\label{sec:hexagonal_inner_purpose}

The hexagon is a visual metaphor, not a geometric requirement. Cockburn has said it was simply a shape that was not yet associated with anything in software architecture. It reminds us that the core should be reachable from many directions, or input devices, without being rewritten for each one. Whether requests come from a GUI, a Web API, a CLI application, a batch job, or an automated test, they should trigger the same use-case behavior.

\section{Step 1: Start with Application Logic}
\label{sec:hexagonal_step1}

At the beginning, we model only application behavior, for example in \texttt{ConfirmKaijuDetectionService} or \texttt{TrackKaijuUseCase}. These classes are responsible for business logic, i.e. executing a use case and enforcing its rules. Commonly such a class is suffixed "Service", or "UseCase", or "Handler". 

At this stage we do not need to commit to \texttt{Spring}, \texttt{ASP.NET}, \texttt{PostgreSQL}, or a message broker. We intentionally start here because architecture decisions become clearer when behavior is visible first.

For the rest of the chapter, we use one running Java example: \texttt{Confirm First Kaiju Detection}. This gives us one consistent slice from driving adapter to core behavior to driven adapters.

\begin{figure}[htbp]
    \centering
    \begin{tikzpicture}[
    core/.style={regular polygon, regular polygon sides=6, shape border rotate=30, draw=blue!70, fill=blue!8, very thick, minimum size=5.8cm, align=center},
    box/.style={rectangle, draw=black!70, fill=white, rounded corners, minimum width=2.7cm, minimum height=0.75cm, align=center, font=\small}
]
    \node[core] (core) {};
    \node[box, yshift=0.75cm] at (core.center) {Use Cases};
    \node[box, yshift=-0.15cm] at (core.center) {Application Services};
    \node[box, yshift=-1.05cm] at (core.center) {Policies and Rules};
\end{tikzpicture}

    \caption{Initial core structure focused on use cases and policies.}
    \label{fig:hexagonal_core_structure}
\end{figure}

The content of the application core has several names. Often, classes in here are called "application services" or "use cases". And are therefore suffixes with "Service" or "UseCase". The responsibility of these classes is to execute a use case and enforce its rules. For example, \texttt{ConfirmKaijuDetectionService} is a service that confirms a Kaiju detection. The class will then expose a method for this particular use case, and implement the logic for it.

Logic include many things:

\begin{itemize}
    \item Validating input data,
    \item Applying business rules,
    \item Coordinating across output ports,
    \item and persisting the results.
\end{itemize}

As such, the application core is the heart of the application, and over time, it will grow to contain the majority of the application's logic, placed in these service classes.

\subsection{What We Delay on Purpose}
\label{sec:hexagonal_delay_decisions}

We deliberately delay technology choices. This gives us two benefits:
\begin{itemize}
    \item We can validate business behavior early without waiting for full infrastructure setup.
    \item We avoid coupling use-case logic to framework types and transport protocols.
\end{itemize}

\subsection{Java Skeleton of the Core}
\label{sec:hexagonal_java_core_skeleton}

The initial core can be expressed as a use-case-focused Java API and service implementation.

First we need the "entry" interface, i.e. the port that will be used by the driving adapters to invoke the use case. It declares a piece of behaviour the core can execute.

\needspace{5\baselineskip}
\begin{ridercode}{Java}
public interface ConfirmKaijuDetectionPort {
    KaijuId confirmDetection(ConfirmKaijuDetectionCommand command);
}
\end{ridercode}

Next, we have the implementation of the port. It is a service class that implements the interface and contains the use-case logic. Here, we confirm a detection by linking a new \texttt{Kaiju} to its origin \texttt{Breach} and persisting it through \texttt{SaveKaijuPort}.

There is also a second output port, \texttt{PublishKaijuConfirmedPort}, implemented by a driven adapter that publishes a mission event. It is similar to the Observer Design Pattern\todoinref{add reference to Observer Pattern}. This event can trigger downstream actions such as command alerts or tracking workflows.

\needspace{16\baselineskip}
\begin{ridercode}{Java}
public final class ConfirmKaijuDetectionService implements ConfirmKaijuDetectionPort {
    private final LoadBreachPort loadBreachPort;
    private final SaveKaijuPort saveKaijuPort;
    private final PublishKaijuConfirmedPort publishKaijuConfirmedPort;
    
    // Constructor omitted for brevity.

    @Override
    public KaijuId confirmDetection(ConfirmKaijuDetectionCommand command) {
        Optional<Breach> breach = loadBreachPort
                .loadById(command.originBreachId())
                .orThrow();

        Kaiju kaiju = Kaiju.confirmed(command.codename(), command.classification(), breach.id());
        saveKaijuPort.save(kaiju);
        publishKaijuConfirmedPort.publishConfirmed(kaiju.id(), breach.id());
        return kaiju.id();
    }
}
\end{ridercode}

\section{Step 2: Add Driving Adapters and Input Ports}
\label{sec:hexagonal_step2}

Once the core behavior exists, we connect ways to invoke it. These are \textit{driving adapters} (also called primary adapters). Their responsibility is to translate external input into calls on input ports, which are interfaces representing use cases.

Obviously, multiple driving adapters can be used to invoke the same use case. For example, we could have a REST controller, a CLI handler, and a GUI controller all invoking the same use case. This is why we need the input port, to abstract the use case from the driving adapter. In the below diagram I have kept it simple and reduced the number of arrows from driving adapters to input ports.

\needspace{16\baselineskip}
\begin{figure}[htbp]
    \centering
    \begin{tikzpicture}[
    core/.style={regular polygon, regular polygon sides=6, shape border rotate=30, draw=blue!70, fill=blue!8, very thick, minimum size=5.5cm, align=center},
    adapter/.style={rectangle, draw=orange!80, fill=orange!12, thick, rounded corners, minimum width=2.3cm, minimum height=0.85cm, align=center, font=\small},
    port/.style={circle, draw=black!70, fill=white, minimum size=0.45cm, inner sep=0pt, font=\tiny},
    lbl/.style={font=\scriptsize, align=center}
]
    \node[core] (core) {Application Core};

    \node[adapter] (rest) at (-5.1, 1.5) {REST Controller};
    \node[adapter] (cli)  at (-5.1, 0.0) {CLI Handler};
    \node[adapter] (gui)  at (-5.1,-1.5) {GUI Controller};

    \node[port] (p1) at (core.150) {};
    \node[port] (p2) at (core.180) {};
    \node[port] (p3) at (core.210) {};
    \node[lbl] at (-2.5, 2.0) {Inbound Port};
    \node[lbl] at (-2.5, -2.0) {Inbound Port};

    \draw[->, thick] (rest.east) -- (p1);
    \draw[->, thick] (cli.east) -- (p1);
    \draw[->, thick] (cli.east) -- (p2);
    \draw[->, thick] (cli.east) -- (p3);
    \draw[->, thick] (gui.east) -- (p3);
\end{tikzpicture}

    \caption{Driving adapters call input ports to invoke application behavior.}
    \label{fig:hexagonal_driving_adapters}
\end{figure}

We can simplify the diagram a bit, by just adding a new hexagonal layer to the left side to represent the driving adapters:

\needspace{16\baselineskip}
\begin{figure}[htbp]
    \centering
    \begin{tikzpicture}[
    % Existing styles
    core/.style={regular polygon, regular polygon sides=6, shape border rotate=30, draw=blue!70, fill=blue!8, very thick, minimum size=5.5cm, align=center},
    port/.style={circle, draw=black!70, fill=white, minimum size=0.45cm, inner sep=0pt, font=\tiny},
    lbl/.style={font=\scriptsize, align=center},
    % New style for the outer boundary
    boundary/.style={draw=blue!90, solid, very thick, fill=teal!30}
]
    % 1. The Outer Half-Hexagon (The Boundary)
    % We use a larger minimum size (e.g., 8cm) to wrap around the core
    \node[regular polygon, regular polygon sides=6, shape border rotate=30, minimum size=8cm] (outer) at (0,0) {};
    
    % Draw only the left half of the outer hexagon
    \draw[boundary] (outer.90) -- (outer.150) -- (outer.210) -- (outer.270) -- cycle;

    % 2. The Application Core
    \node[core] (core) {Application\\Core};

    % 3. Ports and Labels
    \node[port] (p1) at (core.150) {};
    \node[port] (p2) at (core.180) {};
    \node[port] (p3) at (core.210) {};
    \node[port] (p4) at (core.240) {};
    \node[port] (p3) at (core.120) {};
    \node[lbl, rotate=90] at (-1.9, 0.0) {Input Ports};
    
    % Adjusted label position to sit inside the new boundary
    \node[lbl, rotate=90, font=\large] at (-3.2, 0.0) {Driving adapters};

\end{tikzpicture}
    \caption{Driving adapters to the left of the application core.}
    \label{fig:hexagonal_driving_adapters_simplified}
\end{figure}

\subsection{Examples of Driving Adapters}
\label{sec:hexagonal_driving_examples}

Common driving adapters include:
\begin{itemize}
    \item a REST controller translating HTTP requests into use-case calls,
    \item a CLI handler translating command arguments into use-case calls,
    \item a GUI presenter/controller translating user events into use-case calls,
    \item a message consumer translating broker events into use-case calls,
    \item and automated tests that call input ports directly.
\end{itemize}

\subsection{Input Port and Driving Adapter in Java}
\label{sec:hexagonal_java_driving_adapter}

The driving adapter talks only to the input port and maps transport-specific models to use-case models:

\needspace{16\baselineskip}
\begin{ridercode}{Java}
@RestController
@RequestMapping("/kaiju-detections")
public class KaijuController {
    // Dependency injected.
    private final ConfirmKaijuDetectionPort confirmKaijuDetectionPort;

    @PostMapping
    public ResponseEntity<String> confirm(@RequestBody KaijuDetectionRequest request) {
        ConfirmKaijuDetectionCommand command =
                new ConfirmKaijuDetectionCommand(
                        request.codename(),
                        request.classification(),
                        request.originBreachId());
        KaijuId kaijuId = confirmKaijuDetectionPort.confirmDetection(command);
        return ResponseEntity.accepted().body(kaijuId.value());
    }
}
\end{ridercode}

\section{Step 3: Add Driven Adapters and Output Ports}
\label{sec:hexagonal_step3}

The core eventually needs support from external capabilities such as storage, notifications, or third-party APIs. We model these needs as output ports, and implement those ports with \textit{driven adapters} (secondary adapters). This is where dependency inversion becomes concrete: the core depends on interfaces (and \textit{owns} those interfaces), while infrastructure depends on those same interfaces.\autocite{fowler2002patterns}

\begin{figure}[htbp]
    \centering
    \begin{tikzpicture}[
    core/.style={regular polygon, regular polygon sides=6, shape border rotate=30, draw=blue!70, fill=blue!8, very thick, minimum size=5.5cm, align=center},
    adapter/.style={rectangle, draw=orange!80, fill=orange!12, thick, rounded corners, minimum width=2.5cm, minimum height=0.85cm, align=center, font=\small},
    port/.style={circle, draw=black!70, fill=white, minimum size=0.45cm, inner sep=0pt},
    lbl/.style={font=\scriptsize, align=center}
]
    \node[core] (core) {Application Core};

    \node[adapter] (db)    at (5.4, 1.6) {Database Adapter};
    \node[adapter] (email) at (5.4, 0.0) {Email Adapter};
    \node[adapter] (pay)   at (5.4,-1.6) {Payment API Adapter};

    \node[port] (p1) at (core.30) {};
    \node[port] (p2) at (core.0) {};
    \node[port] (p3) at (core.330) {};
    \node[lbl] at (2.5, 2.0) {Output Port};
    \node[lbl] at (2.5, -2.0) {Output Port};

    \draw[<-, thick] (p1) -- (db.west);
    \draw[<-, thick] (p2) -- (email.west);
    \draw[<-, thick] (p3) -- (pay.west);
\end{tikzpicture}

    \caption{Driven adapters implement output ports for infrastructure concerns.}
    \label{fig:hexagonal_driven_adapters}
\end{figure}

Again, we can simplify the diagram a bit, by just adding a new hexagonal layer to the right side to represent the driven adapters:

\needspace{16\baselineskip}
\begin{figure}[htbp]
    \centering
    \begin{tikzpicture}[
    core/.style={regular polygon, regular polygon sides=6, shape border rotate=30, draw=blue!70, fill=blue!8, very thick, minimum size=5.5cm, align=center},
    port/.style={circle, draw=black!70, fill=white, minimum size=0.45cm, inner sep=0pt, font=\tiny},
    lbl/.style={font=\scriptsize, align=center},
    % Boundaries
    driving_bg/.style={draw=blue!90, solid, very thick, fill=teal!30},
    driven_bg/.style={draw=blue!90, solid, very thick, fill=teal!30} % Different color to distinguish
]
    % 1. Reference Outer Hexagon (Invisible)
    \node[regular polygon, regular polygon sides=6, shape border rotate=30, minimum size=8cm] (outer) at (0,0) {};
    
    % 2. Drawing the Two Halves
    % Left Half (Driving)
    \draw[driving_bg] (outer.90) -- (outer.150) -- (outer.210) -- (outer.270) -- cycle;
    % Right Half (Driven)
    \draw[driven_bg] (outer.90) -- (outer.30) -- (outer.330) -- (outer.270) -- cycle;

    % 3. The Application Core
    \node[core] (core) {Application\\Core};

    % 4. Driving Ports (Left)
    \node[port] (lp1) at (core.120) {};
    \node[port] (lp2) at (core.150) {};
    \node[port] (lp3) at (core.180) {};
    \node[port] (lp4) at (core.210) {};
    \node[port] (lp5) at (core.240) {};
    
    % 5. Driven Ports (Right)
    \node[port] (rp1) at (core.60) {};
    \node[port] (rp2) at (core.30) {};
    \node[port] (rp3) at (core.0) {};
    \node[port] (rp4) at (core.330) {};
    \node[port] (rp5) at (core.300) {};

    % Labels
    \node[lbl, rotate=90] at (-1.9, 0.0) {Input Ports};
    \node[lbl, rotate=-90] at (1.9, 0.0) {Output Ports};
    
    % Boundary Labels
    \node[lbl, rotate=90, font=\large\bfseries] at (-3.0, 0.0) {Driving Adapters};
    \node[lbl, rotate=-90, font=\large\bfseries] at (3.0, 0.0) {Driven Adapters};

\end{tikzpicture}
    \caption{Driven adapters to the right of the application core.}
    \label{fig:hexagonal_driven_adapters_simplified}
\end{figure}

And, again, we can update the diagram to include a thin hexagon around the core to represent the ports of the core. Just to clarify that there is a "border" between the core and the adapters, on both sides.

\needspace{16\baselineskip}
\begin{figure}[htbp]
    \centering
    \begin{tikzpicture}[
    % Core style
    core/.style={regular polygon, regular polygon sides=6, shape border rotate=30, draw=blue!70, fill=blue!8, dashed, minimum size=4.5cm, align=center},
    % Ports Layer
    ports_layer/.style={regular polygon, regular polygon sides=6, shape border rotate=30, draw=blue!70, fill=blue!8, very thick, minimum size=5.75cm},
    % Port nodes
    port/.style={circle, draw=black!70, fill=white, minimum size=0.45cm, inner sep=0pt, font=\tiny},
    lbl/.style={font=\scriptsize, align=center},
    % Boundaries
    driving_bg/.style={draw=blue!90, solid, very thick, fill=teal!30},
    driven_bg/.style={draw=blue!90, solid, very thick, fill=teal!30}
]
    % 1. Reference Outer Hexagon (8cm)
    \node[regular polygon, regular polygon sides=6, shape border rotate=30, minimum size=8cm] (outer) at (0,0) {};
    
    % 2. Drawing the Two Boundary Halves
    \draw[driving_bg] (outer.90) -- (outer.150) -- (outer.210) -- (outer.270) -- cycle;
    \draw[driven_bg] (outer.90) -- (outer.30) -- (outer.330) -- (outer.270) -- cycle;

    % 3. The Middle Ports Layer
    \node[ports_layer] (layer) at (0,0) {};
    
    % --- NEW: Vertical line through the Ports Hexagon ---
    \draw[blue!40, very thick] (layer.90) -- (layer.270);

    % 4. The Application Core
    \node[core] (core) {Application\\Core};

    % Labels
    \node[lbl, rotate=90] at (-2.2, 0.0) {Inbound Ports};
    \node[lbl, rotate=-90] at (2.2, 0.0) {Outbound Ports};
    
    % Boundary Labels
    \node[lbl, rotate=90, font=\large\bfseries] at (-3.0, 0.0) {Driving Adapters};
    \node[lbl, rotate=-90, font=\large\bfseries] at (3.0, 0.0) {Driven Adapters};

\end{tikzpicture}
    \caption{Both sides of the hexagon, with the core in the middle.}
    \label{fig:hexagonal_both_sides_simplified}
\end{figure}

Notice how the inner hex is dashed, to indicate a "border" between the core and its ports, but to still indicate that the ports are \textit{inside} the core.

\subsection{Examples of Driven Adapters}
\label{sec:hexagonal_driven_examples}

Typical driven adapters are:
\begin{itemize}
    \item a repository adapter for \texttt{PostgreSQL} or \texttt{MongoDB},
    \item an email adapter for \texttt{SMTP},
    \item a payment adapter for an external payment service API, like Stripe or PayPal,
    \item and a publisher adapter for a message broker, like Kafka or RabbitMQ.
\end{itemize}

\subsection{Output Ports and a Driven Adapter in Java}
\label{sec:hexagonal_java_driven_adapter}

On the driven side, the core defines output ports, and infrastructure implements them.

Let's start with the first output port, \texttt{LoadBreachPort}. It is used to load a Breach from storage. I use \texttt{Optional} \todoinref{add reference, when Optional pattern is written} to indicate that the Breach may not be found.

\needspace{4\baselineskip}
\begin{ridercode}{Java}
public interface LoadBreachPort {
    Optional<Breach> loadById(BreachId breachId);
}
\end{ridercode}

Then the next port to save a Kaiju entity to the storage (database or other storage). I keep the interfaces small here. But sometimes you may want to group multiple related operations into a single interface. For example, you could have an interface which contains both the save and load operations. 

\needspace{4\baselineskip}
\begin{ridercode}{Java}
public interface SaveKaijuPort {
    void save(Kaiju kaiju);
}
\end{ridercode}

And an implementation of the latter port. Here we use \texttt{Spring Data JPA} to save the Kaiju to storage. If you are unfamiliar with \texttt{Spring Data JPA}, you can think of it as a library that provides a way to interact with a database, without explicitly writing SQL queries.

\needspace{8\baselineskip}
\begin{ridercode}{Java}
public final class JpaKaijuRepositoryAdapter implements SaveKaijuPort {
    // Dependency injected.    
    private final SpringDataKaijuRepository repository; 

    @Override
    public void save(Kaiju kaiju) {
        repository.save(KaijuEntity.fromDomain(kaiju));
    }
}
\end{ridercode}

The purpose of this adapter is to hide the database implementation details from the core.

\subsection{Mapping Boundaries in Java}
\label{sec:hexagonal_java_mapping_boundaries}

Mapping code belongs in adapters, because adapters translate between outside models and core models. For simplicity, I use here record classes to represent the commands and requests.

\needspace{8\baselineskip}
\begin{ridercode}{Java}
public record KaijuDetectionRequest(String codename,
                                    KaijuClassification classification,
                                    BreachId originBreachId) {}

public record ConfirmKaijuDetectionCommand(String codename,
                                           KaijuClassification classification,
                                           BreachId originBreachId) {}

ConfirmKaijuDetectionCommand command =
        new ConfirmKaijuDetectionCommand(
                request.codename(),
                request.classification(),
                request.originBreachId());
\end{ridercode}

\needspace{13\baselineskip}
\begin{ridercode}{Java}
@Entity
public class KaijuEntity {
    @Id
    private String id;
    private String codename;

    public static KaijuEntity fromDomain(Kaiju kaiju) {
        KaijuEntity entity = new KaijuEntity();
        entity.id = kaiju.id().value();
        entity.codename = kaiju.codename();
        return entity;
    }
}
\end{ridercode}





\subsection{Dependency Rules}
\label{sec:hexagonal_dependency_rules}

We must remember the dependency rules of this architecture
\begin{itemize}
    \item The core should not depend on the adapters.
    \item The core defines the ports (interfaces) that the adapters must implement.
    \item The adapters should depend on the core.
    \item The adapters should not depend on each other.
\end{itemize}

This means that dependencies should only go from the outside towards to core. \tododiagram{add diagram showing the direction of dependencies. Inner hex and outer hex divided vertically into two parts, with the core in the middle.}
When we then organize the code into packages or modules, we apply these dependency rules by being careful about which packages a class can import.

With this layout, we keep dependency directions explicit:
\begin{itemize}
    \item \texttt{domain} and \texttt{application} do not depend on \texttt{adapters}.
    \item \texttt{adapters.in.*} depend on input ports in \texttt{application.port.in} (i.e. the interfaces providing access to the core). Classes in this package are our driving adapters.
    \item \texttt{application} depends on output ports in \texttt{application.port.out} (i.e. the interfaces that the core defines, to be implemented by the adapters).
    \item \texttt{adapters.out.*} implement output ports and may depend on framework libraries. Classes in this package are our driven adapters.
    \item \texttt{config} performs wiring and can see both core and adapters.
\end{itemize}


\subsection{Requires vs Provides Contracts}
\label{sec:hexagonal_requires_vs_provides}

The most important design choice on the driven side is contract ownership. This is where the hexagonal architecture differs from layered architecture. In hexagonal architecture, the core defines contracts by what the use cases \textit{require}. For example, a core service may depend on \texttt{LoadBreachPort} and \texttt{SaveKaijuPort}, because those capabilities are needed to complete the use case and preserve the Kaiju-origin invariant.

In many layered implementations, interfaces are shaped by what the persistence layer \textit{provides}, for example a repository abstraction aligned to ORM operations. That approach can still be effective, and we can still mock those interfaces. The difference is who drives the abstraction: in hexagonal architecture, required core behavior drives the port definition; in provider-oriented layering, available infrastructure capabilities often drive the interface definition.

This distinction matters during development. With core-owned output ports, we can implement and test use cases before deciding the final database, broker, or external API integration.

\tododiagram{I don't like this diagram}
\begin{figure}[htbp]
    \centering
    \begin{tikzpicture}[
    core/.style={regular polygon, regular polygon sides=6, shape border rotate=30, draw=blue!70, fill=blue!8, very thick, minimum size=3.8cm, align=center, font=\scriptsize},
    port/.style={rectangle, draw=black!70, fill=white, rounded corners, minimum width=2.2cm, minimum height=0.7cm, align=center, font=\scriptsize},
    infra/.style={rectangle, draw=orange!80, fill=orange!12, thick, rounded corners, minimum width=2.3cm, minimum height=0.8cm, align=center, font=\scriptsize},
    title/.style={font=\small}
]
    % Left: hexagonal requires
    \node[title] at (-3.8,2.4) {Hexagonal (Core Defines Requirements)};
    \node[core] (coreL) at (-3.8,0.2) {Application\\Core};
    \node[port] (portL) at (-3.8,-1.8) {\texttt{LoadOrderPort}};
    \node[infra] (infraL) at (-3.8,-3.3) {DB Adapter Implements Port};
    \draw[->, thick] (coreL) -- (portL);
    \draw[->, thick] (infraL) -- (portL);

    % Right: provider-oriented layered tendency
    \node[title] at (3.8,2.4) {Common Layered Tendency (Provider Defines Interface)};
    \node[infra] (infraR) at (3.8,0.2) {Persistence Layer};
    \node[port] (portR) at (3.8,-1.8) {Repository Interface};
    \node[core] (coreR) at (3.8,-3.3) {Application\\Logic Uses It};
    \draw[->, thick] (infraR) -- (portR);
    \draw[->, thick] (coreR) -- (portR);
\end{tikzpicture}

    \caption{Contract ownership comparison: core-defined required ports vs provider-defined interfaces.}
    \label{fig:hexagonal_requires_vs_provides}
\end{figure}

\section{Step 4: Evolving the Core with a Domain Model}
\label{sec:hexagonal_step4}

In some projects, the first hexagonal version uses lightweight application services and simple data structures. That is a valid starting point. As rules become richer, we often add an explicit Domain Model to keep invariants and behavior close to business concepts.

Cockburn's original model did not include an explicit domain model, but if you were to study this architecture online, you will find that it is usually included.



\subsection{Why the Domain Model May Come Later}
\label{sec:hexagonal_domain_later}

In the original hexagonal architecture, there was no domain model. You \textit{can} write software without a domain model of classes and objects. You might go with a simple \textit{data model}, meaning a collection of classes that represent the data you are working with. I realize some people \textit{do} consider this a domain model, but it is an \textit{anemic domain model}, meaning only data but no behavior. You can also go even more primitive and just use simple data structures, but I really don't think anyone wants that.

About the same time as hexagonal architecture was introduced, the Domain-Driven Design (DDD) paradigm was introduced by Eric Evans \autocite{evans2003ddd}. And so, shortly after, the original hexagonal architecture diagram was extended with a domain model.

Adding a domain model later is not a failure of architecture; it is often a sign of responsible evolution. We first establish boundaries and interaction patterns, then deepen the model when complexity justifies it.

We can illustrate the addition of the Domain Model by introducing a new hexagon, inside the core. This now means that the Service classes, in Application Core, can act upon the Domain Model, and that the Domain Model has no dependencies on any outer hexagon.

Figure~\ref{fig:hexagonal_with_domain} shows the updated diagram:

\needspace{16\baselineskip}
\begin{figure}[htbp]
    \centering
    \resizebox{0.6\textwidth}{!}{%  % Adjust 0.8 to 1.0 for full width
        \begin{tikzpicture}
    % 1. Outer Boundary
    \fill[teal!30] (90:4) -- (150:4) -- (210:4) -- (270:4) -- cycle;
    \fill[teal!30] (90:4) -- (30:4) -- (330:4) -- (270:4) -- cycle;
    \draw[blue!90, very thick] (90:4) -- (150:4) -- (210:4) -- (270:4) -- (330:4) -- (30:4) -- cycle;
    \draw[blue!90, very thick] (90:4) -- (270:4);

    % 2. Ports Layer
    \fill[blue!8] (30:2.87) -- (90:2.87) -- (150:2.87) -- (210:2.87) -- (270:2.87) -- (330:2.87) -- cycle;
    \draw[blue!70, very thick] (30:2.87) -- (90:2.87) -- (150:2.87) -- (210:2.87) -- (270:2.87) -- (330:2.87) -- cycle;

    % 3. Application Core
    \draw[blue!70, dashed] (30:2.25) -- (90:2.25) -- (150:2.25) -- (210:2.25) -- (270:2.25) -- (330:2.25) -- cycle;

    % 4. Domain Model
    \fill[blue!30] (30:0.9) -- (90:0.9) -- (150:0.9) -- (210:0.9) -- (270:0.9) -- (330:0.9) -- cycle;
    \draw[blue!90, thick] (30:0.9) -- (90:0.9) -- (150:0.9) -- (210:0.9) -- (270:0.9) -- (330:0.9) -- cycle;

    % 5. Text Labels
    \node at (0,0) [font=\footnotesize\bfseries, align=center, text width=2cm] {Domain\\Model};
    \node at (0,1.2) [blue!80, font=\tiny\bfseries] {Application Core};
    
    \node[rotate=90] at (-2.2, 0) [font=\tiny] {Inbound Ports};
    \node[rotate=-90] at (2.2, 0) [font=\tiny] {Outbound Ports};
    
    \node[rotate=90] at (-3.0, 0) [font=\small\bfseries] {Driving Adapters};
    \node[rotate=-90] at (3.0, 0) [font=\small\bfseries] {Driven Adapters};

    % 6. Dependencies Arrow
    % Points from the bottom left toward the Domain Model boundary
    \draw[-{Stealth[length=5mm, width=3mm]}, very thick, black!80] (-3.5, -3.5) -- (-0.4, -0.7);
    \node at (-2.1, -1.9) [
        rotate=42, fill=white,          % This creates the "box"
        inner sep=1.0pt,     % This controls how much white space is around the letters
        rounded corners=1pt
    ] {Dependencies};

\end{tikzpicture}
    }
    \caption{The domain model inside the core.}
    \label{fig:hexagonal_with_domain}
\end{figure}

For the sake of clarity, I have added a nice arrow to indicate that the dependencies point inward. Outer hexes (layers) may depend on inner hexes, e.g. by importing and using classes or functions from the inner hex (package/module).

Another way to diagram the final architecture is the more common way of using boxes for each layer/part. This is shown in Figure~\ref{fig:hexagonal_as_boxes}:

\needspace{16\baselineskip}
\begin{figure}[htbp]
    \centering
    \resizebox{0.6\textwidth}{!}{%  % Adjust 0.8 to 1.0 for full width
        \usetikzlibrary{arrows.meta, positioning}

\begin{tikzpicture}[
    node distance=1.5cm and 1cm,
    % Styles
    box/.style={draw=blue!90, fill=blue!8, very thick, align=center, minimum width=3.5cm, minimum height=1.2cm},
    port_box/.style={draw=blue!70, fill=blue!15, dashed, font=\tiny\bfseries, minimum height=0.5cm, minimum width=2.4cm},
    adapter_box/.style={draw=blue!90, fill=teal!30, very thick, align=center, minimum width=3.5cm, minimum height=1cm, font=\small\bfseries},
    domain_box/.style={draw=blue!90, fill=blue!30, very thick, align=center, minimum width=3.5cm, minimum height=1cm, font=\small\bfseries},
    arrow/.style={-{Stealth[length=3mm, width=2mm]}, thick, gray!80}
]

    % 1. APPLICATION CORE
    \node[box, minimum height=1.5cm, minimum width=5cm] (app_core) at (0,0) {};
    % Nudged label slightly lower so it doesn't crowd the ports
    \node[anchor=south, yshift=2pt] at (app_core.south) {\scriptsize Application Core};

    % 2. PORTS
    \node[port_box, anchor=north west] (in_port) at (app_core.north west) {Input Ports};
    \node[port_box, anchor=north east] (out_port) at (app_core.north east) {Output Ports};

    % 3. DOMAIN MODEL - Changed gap from 1.5cm (default) to 0.7cm
    \node[domain_box, below=0.7cm of app_core] (domain) {Domain\\Model};

    % 4. ADAPTERS
    \node[adapter_box, above left=1.0cm and -2.4cm of in_port] (driving) {Driving Adapters};
    \node[adapter_box, above right=1.0cm and -2.4cm of out_port] (driven) {Driven Adapters};

    % 5. ARROWS
    \draw[arrow] (driving.south) -- (in_port.north);
    \draw[arrow] (driven.south) -- (out_port.north);

    % Core to Domain
    \draw[arrow] (app_core.south) -- (domain.north);

\end{tikzpicture}
    }
    \caption{The final architecture as boxes.}
    \label{fig:hexagonal_as_boxes}
\end{figure}


\section{Package Layout}
\label{sec:hexagonal_package_layout}

An important aspect of \textit{architecture} is the organization of the code into packages and modules, along with the dependencies between them. Most architectural styles will have a recommended superficial package layout, so let us explore how we can organize an application using hexagonal architecture.

A package convention gives us a practical guardrail. We can often organize our code by architectural role, so that core code remains stable while adapters can change. There are many ways to do this, I will give an overall recommendation, and then two suggestions on how to organize the aplication hex: layers and feature-folders. 

\subsection{Overall structure}

We generally have three parts to a hexagonal architecture: the core, the driving (inbound) adapters, and the driven (outbound) adapters. So, that's where we start: create three packages (or modules).

\needspace{5\baselineskip}
\begin{alltt}\small
    \foldericon\ src/
    |-- \foldericon\ application/
    |-- \foldericon\ inboundadapters/
    `-- \foldericon\ outboundadapters/
\end{alltt}

Many applications, across many architectural styles, will actually have a similar package structure, perhaps just differently named. Sometimes the application is called "core", sometimes "services". The same responsibility has been named different things over the years.

Next, we also need the domain. We can place it inside the application package, to represent that the domain hex is inside the application hex. Or, we can put it in a separate package. Below, I will show the latter:

\needspace{6\baselineskip}
\begin{alltt}\small
    \foldericon\ src/
    |-- \foldericon\ application/
    |-- \foldericon\ domain/
    |-- \foldericon\ inboundadapters/
    `-- \foldericon\ outboundadapters/
\end{alltt}

Then, we have to further organize the application hex, or package. As mentioned earlier, I will show two approaches: layers and feature-folders.

\subsection{Application Layers}

Let's start with layers. The application hex consists of service classes (the classes performing the use cases), inbound ports (the interfaces exposing the use cases to the outside world), and outbound ports (the interfaces defining the capabilities the core needs from the outside world). If we go with the layered approach, we divide files by their responsibility, or category. That means three sub-packages:

\needspace{10\baselineskip}
\begin{alltt}\small
    \foldericon\ src/
    |-- \foldericon\ application/
    |   |-- \foldericon\ inboundports/
    |   |-- \foldericon\ outboundports/
    |   `-- \foldericon\ services/
    |-- \foldericon\ domain/
    |-- \foldericon\ inboundadapters/
    `-- \foldericon\ outboundadapters/
\end{alltt}

This is the basics. Your application may need other kinds of classes, which do not fit into one of the three categories. 

I will not go into too details about the adapters here. In your \texttt{inboundadapters} package, you will have classes that let's the outside world interact with the application. This could be a Web API controller, a CLI handler, a GUI controller, a message consumer, or something else. Let's just say that our application exposes a REST API, and also has a GUI. I will make two sub-packages in the \texttt{inboundadapters} package: \texttt{web} and \texttt{gui}.

\needspace{10\baselineskip}
\begin{alltt}\small
    \foldericon\ src/
    |-- \foldericon\ application/
    |   |-- \foldericon\ inboundports/
    |   |-- \foldericon\ outboundports/
    |   `-- \foldericon\ services/
    |-- \foldericon\ domain/
    |-- \foldericon\ inboundadapters/
    |   |-- \foldericon\ web/
    |   `-- \foldericon\ gui/
    `-- \foldericon\ outboundadapters/
\end{alltt}

What about outbound adapters? Again, we have many categories, try to organize your code by your needs. For example, you may have a package for \texttt{persistence}, \texttt{events}, \texttt{notifications}, \texttt{messaging}, \texttt{email}, \texttt{logging}, \texttt{sensordrivers}, \texttt{externalAPIs}, etc. Let's just add a sub-package for \texttt{persistence} and \texttt{email}:

\needspace{12\baselineskip}
\begin{alltt}\small
    \foldericon\ src/
    |-- \foldericon\ application/
    |   |-- \foldericon\ inboundports/
    |   |-- \foldericon\ outboundports/
    |   `-- \foldericon\ services/
    |-- \foldericon\ domain/
    |-- \foldericon\ inboundadapters/
    |   |-- \foldericon\ gui/
    |   `-- \foldericon\ web/
    `-- \foldericon\ outboundadapters/
        |-- \foldericon\ email/
        `-- \foldericon\ persistence/
\end{alltt}

Further subdivision of your application is of course possible, and as the capabilities of the application grow, you will likely need to add more sub-packages. Here, you will have to decide what is the best way to organize your code.

\subsection{Feature Folders}

Another way to organize the application is to use feature folders. In the previous approach, a service class in the \texttt{services} package will expose an inboundport, located in the \texttt{inboundports} package. And it may depend on a number of outbound ports, located in the \texttt{outboundports} package. That means, classes which work closely together to implement a feature, are spread across different packages, and navigating between those files may become cumbersome. That, at least, is the concern, which leads to the feature folder approach. So, instead, we organize the code by feature, or use case, a package per. And the service class, with its inbound and outbound ports, will be located in the feature package.

Here is an example, I have just updated the \texttt{application} package to use feature folders:

\needspace{12\baselineskip}
\begin{alltt}\small
    \foldericon\ src/
    |-- \foldericon\ application/
    |   `-- \foldericon\ features/
    |       |-- \foldericon\ confirmkaijudetection/
    |       |   |-- \fileicon\ ConfirmKaijuDetectionInboundPort.java
    |       |   |-- \fileicon\ ConfirmKaijuDetectionService.java
    |       |   |-- \fileicon\ KaijuConfirmedEvent.java
    |       |   `-- \fileicon\ ConfirmKaijuDetectionOutboundPort.java
    |       |-- \foldericon\ registerkaiju/
    |       `-- \foldericon\ triangulatekaiju/
    |-- \foldericon\ domain/
    |-- \foldericon\ inboundadapters/
    |   |-- \foldericon\ gui/
    |   `-- \foldericon\ web/
    `-- \foldericon\ outboundadapters/
        |-- \foldericon\ email/
        `-- \foldericon\ persistence/
\end{alltt}

Notice how each feature package clearly describes the feature through its name. It then contains whatever classes and interfaces are needed to implement the feature. The adapters, of course, are still placed outside. Just pretend the other feature packages also contains relevant files.

You might similarly organize your adapter layers by feature, so in the \texttt{outboundadapters} package, you will have sub-packages \texttt{confirmkaijudetection}, \texttt{registerkaiju}, and \texttt{triangulatekaiju}. And the same for the \texttt{inboundadapters} package.

\subsection{Layers vs Feature Folders}

So, which approach should you choose? Both are valid, and both have their own advantages and disadvantages.



% ------------------------------------------------------------

\section{Testing Across the Hexagon Boundary}
\label{sec:hexagonal_testing}

The separation of concerns gives us a practical testing strategy:\autocite{huttunen2026hexagonal}
\begin{itemize}
    \item test application behavior through input ports without real infrastructure,
    \item test adapters against their technologies in focused integration tests,
    \item and keep end-to-end tests for wiring and critical user paths.
\end{itemize}

In short, we can focus on unit testing the behaviour of the core, by mocking \todoinref{add reference to mocking} the driven adapters.



\subsection{Why Ports Improve Testability}
\label{sec:hexagonal_why_ports_testability}

Ports improve testability because they create explicit seams \todoinref{add reference to seams} around behavior. An input port gives us a stable way to invoke a use case, while output ports let us replace infrastructure with in-memory test doubles \todoinref{add reference to test doubles}. 
In practice, we execute a use case through an input port and provide fake implementations such as \texttt{LoadBreachPort} or \texttt{PublishKaijuConfirmedPort}. This keeps most tests fast and deterministic, because they verify business rules without requiring a running database, broker, or HTTP stack.

This is not exclusive to hexagonal architecture. A disciplined layered architecture can also isolate dependencies behind interfaces and achieve good tests. The practical difference is that hexagonal architecture makes these seams a first-class design element around the core, so isolated tests become the default rather than an optional refactoring.

\subsection{Testing Slices with Java}
\label{sec:hexagonal_java_testing_slices}

In Java, we can keep most behavior tests focused on the core. Here is an example of testing a use case defined in a service class. We use a fake implementation of the repository and the event publisher. Both fakes implement the the necessary interface.

\needspace{30\baselineskip}
\begin{ridercode}{Java}
@Test
void confirmsDetectionAndPublishesEvent() {
    FakeLoadBreachPort loadPort = new FakeLoadBreachPort();
    FakeSaveKaijuPort savePort = new FakeSaveKaijuPort();
    FakePublishKaijuConfirmedPort publishPort = 
                        new FakePublishKaijuConfirmedPort();

    ConfirmKaijuDetectionService service = 
                    new ConfirmKaijuDetectionService(
                        loadPort,
                        savePort,
                        publishPort
                    );

    ConfirmKaijuDetectionCommand command = 
                    new ConfirmKaijuDetectionCommand(
                        "Leatherback",
                        KaijuClassification.CATEGORY_III,
                        new BreachId("B-01")
                    );
            
    KaijuId id = service.confirmDetection(command);

    assertTrue(savePort.savedKaijuIds().contains(id));

    assertTrue(publishPort
                .publishedKaijuIds()
                .contains(command.originBreachId().value())
                );
}
\end{ridercode}


\section{Practical Trade-offs}
\label{sec:hexagonal_tradeoffs}

Hexagonal architecture improves replaceability of the adapters and testability, but it introduces explicit interfaces and mapping code. For small CRUD-heavy systems, that can feel heavy. For systems with changing workflows, multiple interfaces, or expensive integration points, the added structure usually pays for itself over time.

\subsection{Hexagonal and Layered: A Preview}
\label{sec:hexagonal_layered_preview}

Hexagonal and layered architecture are not enemies; both can be implemented well or poorly. A layered design with clear interfaces and dependency inversion can achieve strong testability. The practical difference we emphasize is orientation: hexagonal architecture centers on inside/outside boundaries and core-owned required ports, while layered architecture commonly centers on presentation, application, and persistence layers.


\subsection{When Migration Is Worth It}
\label{sec:hexagonal_migration_when}

Migration from layered to hexagonal is most useful when integration change is frequent, tests are dominated by slow end-to-end flows, or business rules are spread across controllers and repositories. If those symptoms are not present, a disciplined layered approach may remain the better economic choice.

\subsection{Incremental Migration Steps}
\label{sec:hexagonal_migration_steps}

We can migrate incrementally instead of rewriting everything:
\begin{enumerate}
    \item select one high-value use case (for example \texttt{ConfirmKaijuDetection}),
    \item define an input port and required output ports for that use case,
    \item move business logic into an application service behind the input port,
    \item wrap current repository or API integrations as driven adapters,
    \item add fast core tests that use fake output ports,
    \item then repeat use case by use case.
\end{enumerate}

What we essentially do here is create a port owned by the core, and then implement it with a driven adapter.

\subsection{Common Implementation Mistakes}
\label{sec:hexagonal_common_mistakes}

Even with a good structure, implementations can drift. Typical mistakes include:
\begin{itemize}
    \item placing framework annotations in core domain models and use cases,
    \item exposing \texttt{JPA} entities directly across the core boundary,
    \item defining repository interfaces around provider capabilities rather than use-case needs,
    \item and embedding HTTP or SQL concerns in application services.
\end{itemize}